\documentclass[12pt]{article}
\usepackage[T1]{fontenc}
\usepackage[utf8]{inputenc}
\usepackage{lmodern}
\usepackage{textcomp}
\usepackage{enumitem}
\usepackage{geometry}
\usepackage{graphicx}
\usepackage{amsmath, amssymb}
\geometry{margin=1in}
\begin{document}
\begin{center}
Astronomy - Graduate Level Assessment 10 \
Total Marks: 40 \
Answer all questions.
\end{center}

\section*{Multiple Choice (10 marks)}
\begin{enumerate}[label=\arabic*.]
\item What Mars mission will be landing on May 25 2008 and will dig a trench into (hopefully) ice-rich soil?
\begin{enumerate}[label=(\alph*)]
    \item ExoMars
    \item Mars Exploration Rovers
    \item Mars Science Laboratory
    \item Phoenix Mars Lander
\end{enumerate}
\item The term Schwarzschild radius usually describes properties of ...
\begin{enumerate}[label=(\alph*)]
    \item red dwarfs.
    \item pulsars.
    \item black holes.
    \item galaxies.
\end{enumerate}
\item How do astronomers think Jupiter generates its internal heat?
\begin{enumerate}[label=(\alph*)]
    \item nuclear fusion in the core
    \item by contracting changing gravitational potential energy into thermal energy
    \item internal friction due to its high rotation rate
    \item chemical processes
\end{enumerate}
\item The nebular theory of the formation of the solar system successfully predicts all but one of the following. Which one does the theory not predict?
\begin{enumerate}[label=(\alph*)]
    \item Planets orbit around the Sun in nearly circular orbits in a flattened disk.
    \item the equal number of terrestrial and jovian planets
    \item the craters on the Moon
    \item asteroids Kuiper-belt comets and the Oort cloud
\end{enumerate}
\item The famous Drake equation attempts to answer the following question:
\begin{enumerate}[label=(\alph*)]
    \item Will the Sun become a black hole?
    \item Is the universe infinitely large?
    \item How old is the visible universe?
    \item Are we alone in the universe?
\end{enumerate}
\end{enumerate}

\section*{True / False (10 marks)}
\textit{Answer True or False}
\begin{enumerate}[label=\arabic*.]
\setcounter{enumi}{5}
\item True or False: When a blackbody is cooled to half its original temperature, the power emitted decreases to 1/16 of its original value and the peak wavelength doubles.
\item True or False: Knowing an object's actual brightness and its apparent brightness allows for the estimation of its distance from the observer.
\item True or False: Planet X is tidally locked, resulting in the absence of a solar day as it does not complete a rotation relative to its sun.
\item True or False: Bolometric luminosity refers specifically to the luminosity integrated over visible wavelengths only.
\item True or False: Asteroids emit a significant amount of their own radiation, which makes them stand out in sky surveys.
\end{enumerate}

\section*{Long Form Response (20 marks)}
\textit{Answer each question with a clear and complete explanation.}
\begin{enumerate}[label=\arabic*.]
\setcounter{enumi}{10}
\item Discuss the major hypotheses regarding the origin of the Moon, and critically analyze why the tidal forces leading to a split from Earth is not a widely accepted explanation.
\item Discuss the significance of Betelgeuse within the Orion constellation, including its characteristics, cultural impact, and role in stellar evolution studies.
\end{enumerate}

\vfill
\noindent\textit{End of Paper}
\end{document}